\chapter*{Big Picture}
\addcontentsline{toc}{chapter}{Big Picture}
\label{front:big-picture}

The goal of this book is to understand one object from three directions.  The
object is a genus-one complex group.  It first appears analytically as a quotient
\[
  T=\C/\Lambda,
\]
where addition is inherited from addition in \(\C\).  It then appears
geometrically as a Jacobian,
\[
  \Pic^0(E),
\]
where the group law is encoded by degree-zero divisor classes.  Finally it
appears algebraically as the complex points of a smooth plane cubic,
\[
  E(\C):\quad Y^2Z=4X^3-g_2XZ^2-g_3Z^3.
\]

\begin{center}
\resizebox{0.82\linewidth}{!}{\classicTrinityDiagram}
\end{center}

\noindent
The central theorem says that these are not merely related examples:
\[
  \C/\Lambda \cong \Pic^0(E) \cong E(\C)
\]
as complex Lie groups.  This phrase is important: the maps are both
biholomorphic and compatible with the group laws.  Thus the statement is
stronger than an isomorphism of abstract abelian groups; it also preserves the
complex structure.  The most visible consequence is the chord--tangent law.  On
the torus, addition is simply
\[
  [u]+[v]=[u+v].
\]
On the cubic, the same addition is disguised as a geometric recipe: draw the line
through two points, take the third intersection with the cubic, and reflect it.

\section*{What Must Be Understood}

\begin{enumerate}
\item \textbf{Why a torus is a Riemann surface.}
  A lattice quotient \(\C/\Lambda\) is compact, genus one, and still carries a
  holomorphic group law.

\item \textbf{Why elliptic functions are constrained.}
  Meromorphic periodic functions cannot have arbitrary zeros and poles.  The
  residue theorem and the argument principle force balancing laws.

\item \textbf{Why the Weierstrass function creates a cubic.}
  The functions \(\wp\) and \(\wp'\) satisfy
  \[
    \wp'^2=4\wp^3-g_2\wp-g_3,
  \]
  so \(z\mapsto[\wp(z):\wp'(z):1]\) sends the torus into a plane cubic.

\item \textbf{Why divisors produce the Jacobian.}
  Principal divisors have degree zero, and Abel's theorem says that on a torus a
  degree-zero divisor is principal exactly when its Abel sum vanishes.

\item \textbf{Why Riemann--Roch gives the correct coordinates.}
  On a genus-one surface, \(\ell(nO)=n\).  Thus \(1,\wp,\wp'\) are exactly the
  three functions needed to map the torus into \(\CP^2\).

\item \textbf{Why line intersection becomes addition.}
  A line meeting the cubic in \(P,Q,R\) cuts out a principal divisor
  \([P]+[Q]+[R]-3[O]\).  Abel's theorem converts this into the group equation
  \(P+Q+R=O\), so the sum of \(P\) and \(Q\) is the reflection of \(R\).
\end{enumerate}

\begin{keyidea}{The route through the book}
The chapters move from the exam tools to the torus, from the torus to elliptic
functions, from elliptic functions to the cubic, and from divisors to the
Jacobian.  The final synthesis is that the analytic, geometric, and algebraic
constructions describe the same group.
\end{keyidea}
